\documentclass[11pt]{article}
\usepackage[T1]{fontenc}
\usepackage{lmodern}
\usepackage{amsmath, amssymb, amsthm}
\usepackage[margin=1in]{geometry}
\usepackage{hyperref}
\usepackage[dvipsnames]{xcolor}
\usepackage{tcolorbox}
\tcbuselibrary{skins, breakable}

% ── tcolorbox colour palette (matches Modular_UpdateDerivations style) ──────
\newtcolorbox{derivbox}[1][]{colback=gray!8, colframe=gray!50,
  fonttitle=\bfseries, title=#1, breakable}
\newtcolorbox{ebnmbox}[1][]{colback=blue!5, colframe=blue!40,
  fonttitle=\bfseries, title=#1, breakable}
\newtcolorbox{verifybox}[1][]{colback=green!5, colframe=green!40,
  fonttitle=\bfseries, title=#1, breakable}
\newtcolorbox{correctionbox}[1][]{colback=red!5, colframe=red!40,
  fonttitle=\bfseries, title=#1, breakable}

\newcommand{\E}{\mathbb{E}}
\newcommand{\Var}{\mathrm{Var}}
\newcommand{\KL}{\mathrm{KL}}
\newcommand{\ELBO}{\mathcal{L}}
\newcommand{\N}{\mathcal{N}}
\newcommand{\diag}{\mathrm{diag}}

\title{\textbf{D1: Simplified $q(\mathbf{L})$ Update Under $\eta = \mathbf{Y}\mathbf{F}\tilde{\boldsymbol{\beta}}$}\\[4pt]
\large Cluster B (Cox-on-YF Reformulation) — Derivation D1\\[2pt]
\normalsize Cross-reference: \texttt{derivations/Modular\_UpdateDerivations/} q($L$) section}
\author{Andrew Walther}
\date{2026-04-29}

\begin{document}
\maketitle

\tableofcontents
\newpage

% ============================================================
\section{Context and Goal}
% ============================================================

In the original SSBMF model, the Cox proportional-hazards linear predictor is
\begin{equation}
  \eta_i = \mathbf{l}_i^\top \boldsymbol{\beta}, \qquad
  \text{where } \mathbf{l}_i = \text{row } i \text{ of the latent loading matrix } \mathbf{L}.
\end{equation}
This makes $\mathbf{L}$ appear in \emph{both} the genomics likelihood ($\mathbf{Y} \approx \mathbf{L}\mathbf{F}^\top$)
and the Cox likelihood (through $\eta_i$), so the variational update for $q(\mathbf{L})$ must balance
two information sources (``dual-source EBNM'').

\bigskip
\textbf{Cluster B reformulation.} Replace the linear predictor with
\begin{equation}
  \eta_i = (\mathbf{y}_i \mathbf{F})\,\tilde{\boldsymbol{\beta}} = \sum_k \left(\sum_j y_{ij} f_{jk}\right)\tilde{\beta}_k,
  \label{eq:new-eta}
\end{equation}
where $\tilde{\boldsymbol{\beta}}$ is the CAVI variational parameter and
$\mathbf{Z}_F = \mathbf{Y}\mathbf{F} \in \mathbb{R}^{n \times K}$ are
\emph{observed} projection scores (held fixed per outer iteration).

\bigskip
\textbf{Goal of D1.} Show that under~\eqref{eq:new-eta}, $\mathbf{L}$ no longer
appears in the Cox likelihood term, so the ELBO's dependence on $q(\mathbf{L})$
reduces to the genomics likelihood alone, yielding a pure-genomics EBNM update.

% ============================================================
\section{ELBO Decomposition Under the New Model}
% ============================================================

The full ELBO is
\begin{equation}
  \ELBO = \underbrace{\E_q[\log p(\mathbf{Y} \mid \mathbf{L}, \mathbf{F}, \boldsymbol{\tau})]}_{\text{genomics}}
         + \underbrace{\E_q[\log p(\mathbf{t}, \boldsymbol{\delta} \mid \mathbf{Y}, \mathbf{F}, \tilde{\boldsymbol{\beta}})]}_{\text{survival (new)}}
         - \KL[q(\mathbf{L}) \| p(\mathbf{L})]
         - \KL[q(\mathbf{F}) \| p(\mathbf{F})]
         - \KL[q(\tilde{\boldsymbol{\beta}}) \| p(\tilde{\boldsymbol{\beta}})]
         - \KL[q(\boldsymbol{\tau}) \| p(\boldsymbol{\tau})].
\end{equation}

\begin{derivbox}[Key observation]
Under~\eqref{eq:new-eta}, the survival term
$\E_q[\log p(\mathbf{t}, \boldsymbol{\delta} \mid \mathbf{Y}, \mathbf{F}, \tilde{\boldsymbol{\beta}})]$
depends on $\mathbf{Z}_F = \mathbf{Y}\mathbf{F}$ and $\tilde{\boldsymbol{\beta}}$, but
\textbf{not on $\mathbf{L}$}. The matrix $\mathbf{L}$ appears only in the genomics likelihood.
\end{derivbox}

\subsection{Survival Term Does Not Involve $\mathbf{L}$}

The Cox partial likelihood (and its Taylor-linearised surrogate used in CAVI) depends on
$\eta_i = \sum_k (z_F)_{ik}\,\tilde{\beta}_k$ where $(z_F)_{ik} = \sum_j y_{ij} f_{jk}$.
The terms $(z_F)_{ik}$ are functions of the \emph{observed} $\mathbf{Y}$ and the \emph{factor} $\mathbf{F}$;
they do not involve the loadings $\mathbf{L}$.

Therefore
\begin{equation}
  \frac{\partial}{\partial q(\mathbf{L})} \E_q[\log p(\mathbf{t}, \boldsymbol{\delta} \mid \mathbf{Y}, \mathbf{F}, \tilde{\boldsymbol{\beta}})] = 0.
\end{equation}

\subsection{Genomics Likelihood Depends on $\mathbf{L}$}

The genomics model is $\mathbf{Y} \approx \mathbf{L}\mathbf{F}^\top + \mathbf{E}$ where
$e_{ij} \overset{\text{iid}}{\sim} \N(0, \tau_j^{-1})$. The expected log-likelihood is
\begin{equation}
  \E_q[\log p(\mathbf{Y} \mid \mathbf{L}, \mathbf{F}, \boldsymbol{\tau})]
  = -\frac{1}{2}\sum_{i,j} \tau_j \,\E_q[(y_{ij} - \mathbf{l}_i^\top \mathbf{f}_j)^2].
  \label{eq:genomics-lik}
\end{equation}
This depends on $q(\mathbf{L})$ through $\E_q[\mathbf{l}_i]$ and $\E_q[l_{ik}^2]$.

% ============================================================
\section{Factor-wise ELBO for q($l_k$)}
% ============================================================

In Gauss-Seidel CAVI, we update $q(\mathbf{l}_k)$ holding all other factors fixed.
Define the partial residual
\begin{equation}
  R_{k} = \mathbf{Y} - \sum_{k' \ne k} \E_q[\mathbf{l}_{k'}]\,\E_q[\mathbf{f}_{k'}]^\top
        \;\in\; \mathbb{R}^{n \times p}.
\end{equation}

Isolating factor $k$ in~\eqref{eq:genomics-lik}:
\begin{align}
  \E_q[\log p(\mathbf{Y} \mid \mathbf{L}, \mathbf{F}, \boldsymbol{\tau})] \Big|_k
  &= -\frac{1}{2}\sum_{i,j} \tau_j \,\E_q[(r_{k,ij} - l_{ik}\,f_{jk})^2] + \text{const} \nonumber\\
  &= -\frac{1}{2}\sum_{i} \left[
       \E_q[l_{ik}^2]\sum_j \tau_j \E_q[f_{jk}^2]
       - 2\,\E_q[l_{ik}]\sum_j \tau_j \E_q[r_{k,ij}]\E_q[f_{jk}]
     \right] + \text{const}.
  \label{eq:factor-k-lik}
\end{align}

\begin{verifybox}[Survival terms: absent]
The factor-$k$ contribution to the ELBO that involves $q(l_{ik})$ is given
entirely by~\eqref{eq:factor-k-lik}. There are \textbf{no} $w_i$, $\tilde{\beta}_k$,
$z_{\text{no-}k}$ or other survival terms present, because $\mathbf{L}$ is absent from
the survival likelihood under the Cluster B reformulation.
\end{verifybox}

% ============================================================
\section{Variational Update for $q(l_{ik})$ — Pure-Genomics EBNM}
% ============================================================

Completing the square in~\eqref{eq:factor-k-lik}, the optimal $q(l_{ik})$ is the EBNM
solution with pseudo-observation
\begin{equation}
  x_{L,ik} = \frac{B_{L,ik}}{A_{L,i}}, \qquad
  s_{L,ik} = \frac{1}{\sqrt{A_{L,i}}},
\end{equation}
where
\begin{align}
  A_{L,i} &= \sum_j \tau_j\,\E_q[f_{jk}^2], \label{eq:AL-gen}\\[4pt]
  B_{L,ik} &= \sum_j \tau_j\,\E_q[r_{k,ij}]\,\E_q[f_{jk}]. \label{eq:BL-gen}
\end{align}

Note that $A_{L,i}$ is the \emph{same} for all patients $i$ (it sums over features $j$),
so the subscript $i$ is redundant in practice — written here for clarity.

\begin{ebnmbox}[EBNM call for q($l_k$) under Cluster B]
\textbf{Input:} \;
$x_{L,\cdot k} = B_{L,\cdot k} / A_{L}$ \,(n-vector),\;
$s_{L} = 1/\sqrt{A_{L}}$ \,(scalar, same for all $i$)

\vspace{6pt}
\textbf{Prior family:} point-exponential (or as configured in \texttt{prior\_LF})

\vspace{6pt}
\textbf{Output:} $\E_q[l_{ik}]$, $\E_q[l_{ik}^2]$, \texttt{ebnm\_log\_likelihood}
\end{ebnmbox}

% ============================================================
\section{Comparison to Original q($L$) Update}
% ============================================================

\begin{center}
\renewcommand{\arraystretch}{1.4}
\begin{tabular}{lll}
\hline
\textbf{Quantity} & \textbf{Original ($\eta = \mathbf{L}\boldsymbol{\beta}$)} & \textbf{Cluster B ($\eta = \mathbf{Z}_F\tilde{\boldsymbol{\beta}}$)} \\
\hline
$A_{L,i}$ & $(1-\alpha)\,A_{\text{gen}} + \alpha\,\lambda\,w_i\,\E[\beta_k^2]$ & $A_{\text{gen}} = \sum_j \tau_j\E[f_{jk}^2]$ \\
$B_{L,ik}$ & $(1-\alpha)\,B_{\text{gen},ik} + \alpha\,\lambda\,w_i\,z_{\text{no-}k,i}\,\E[\beta_k]$ & $B_{\text{gen},ik}$ only \\
Survival args & $w$, $\E[\beta_k]$, $\E[\beta_k^2]$, $z_{\text{no-}k}$, $\lambda$, $\alpha$ & \emph{none} \\
\hline
\end{tabular}
\end{center}

\bigskip
The Cluster B update is strictly simpler: $\alpha$ no longer appears (or equivalently,
the $(1-\alpha)$ factor in front of $A_{\text{gen}}$ becomes 1, so we may drop it and
set the floor directly). The parameters $w$, $\tilde{\beta}_k$,
$\E[\tilde{\beta}_k^2]$, $z_{\text{no-}k}$, $\lambda$, and \texttt{normalize\_AB}
are all removed from \texttt{update\_L\_k()}.

% ============================================================
\section{Code Implications}
% ============================================================

\paragraph{\texttt{code/update\_L.R} — \texttt{update\_L\_k()}}
\begin{itemize}
  \item \textbf{Remove arguments:} \texttt{w}, \texttt{EBeta\_k}, \texttt{EBeta2\_k},
        \texttt{z\_no\_k}, \texttt{lambda}, \texttt{normalize\_AB}, \texttt{alpha}
  \item \textbf{Remove:} \texttt{A\_surv} and \texttt{B\_surv} computation blocks
  \item \textbf{Simplify:} \texttt{A\_L <- pmax(A\_gen, A\_floor)},\;
        \texttt{B\_L <- B\_gen}
  \item \textbf{New signature:}
        \texttt{update\_L\_k(Tau, EF\_k, EF2\_k, R\_k, prior\_family, A\_floor)}
\end{itemize}

\paragraph{\texttt{code/update\_L.R} — \texttt{update\_L\_all()}}
Drop \texttt{w}, \texttt{z}, \texttt{EBeta}, \texttt{EBeta2}, \texttt{lambda},
\texttt{normalize\_AB} from the outer wrapper. The Gauss-Seidel $k$-loop structure
is preserved; only the argument threading changes.

% ============================================================
\section{Verification Checklist}
% ============================================================

\begin{verifybox}[D1 Verification]
\begin{enumerate}
  \item $\eta_i = \sum_k (z_F)_{ik}\,\tilde{\beta}_k$ contains no $l_{ik}$. \checkmark
  \item The Cox Taylor expansion — $\log \mathrm{PL} \approx \mathrm{logPL}_0
        + \sum_i u_i\,\delta\eta_i - \tfrac{1}{2}\sum_i w_i\,(\delta\eta_i)^2$
        — involves $\mathbf{Z}_F$ and $\tilde{\boldsymbol{\beta}}$, not $\mathbf{L}$. \checkmark
  \item Factor-wise ELBO for $q(l_k)$ contains only the genomics
        terms~\eqref{eq:AL-gen}--\eqref{eq:BL-gen}. \checkmark
  \item All survival-specific arguments removed from \texttt{update\_L\_k()}. (to implement)
  \item 171/171 tests pass after simplification. (to implement)
\end{enumerate}
\end{verifybox}

\end{document}
